% Lösungen Einheit 2 – Prozentwert berechnen
\einheitenkopf[][2 Prozentwert]{Einheit 2 von 4 · Prozentwert berechnen}

\erg{24}{a) $10\,€$ \quad b) $30\,€$ \quad c) $70\,€$ \quad d) $150\,€$ \quad e) Der Anteil bleibt $50\,\%$, das Ganze wächst; der Teil ist immer die Hälfte des Ganzen.}
\erg{25}{a) $15\,$kg \quad b) $7\,$m \quad c) $7{,}5\,$kg}
\erg{26}{a) $1\,\% = 3\,€$, $7\,\% = 21\,€$ \quad b) $1\,\% = 0{,}5\,$kg, $6\,\% = 3\,$kg}
\erg{27}{a) $10\,\% = 9\,€$, $30\,\% = 27\,€$ \quad b) $10\,\% = 6\,€$, $5\,\% = 3\,€$, $15\,\% = 9\,€$ \quad c) $10\,\% = 14\,€$, $5\,\% = 7\,€$}
\erg{28}{a) $0{,}4 \cdot 30\,€ = 12\,€$ \quad b) $0{,}08 \cdot 250\,€ = 20\,€$ \quad c) $0{,}2 \cdot 4{,}50\,€ = 0{,}90\,€$}
\erg{29}{a) $15\,€$ \quad b) $60\,€ - 15\,€ = 45\,€$ \quad c) Rabatt $35\,€$, neuer Preis $315\,€$ \quad d) $57\,€$ ($0{,}15 \cdot 380\,€$; $323\,€$ ist der neue Preis)}
\erg{30}{a) $60\,€$ \quad b) $1{,}3 \cdot 80\,$cm $= 104\,$cm}
\erg{31}{a) $40\,\%$ ist $0{,}4$, nicht $0{,}04$; richtig $0{,}4 \cdot 60\,€ = 24\,€$ \quad b) Emma hat den neuen Preis ausgerechnet, gefragt ist die Ersparnis; richtig $0{,}2 \cdot 45\,€ = 9\,€$}
\erg{32}{a) $1\,\%$ ist ein Hundertstel; ein Hundertstel einer Größe erhält man, wenn man sie in 100 gleiche Teile teilt. \quad b) Ja; beide Male rechnet man $20 \cdot 50 : 100$, also $10\,€$.}
\erg{33}{a) $1\,\% = 2{,}40\,€$, $19\,\% = 45{,}60\,€$ \quad b) $285{,}60\,€$ \quad c) Der zweite, denn $280\,€ < 285{,}60\,€$.}
